\label{sec:experiments}

\subsection{Dataset and Implementation}

The corpus contains 10{,}741 Vietnamese dishes with structured fields (name, ingredients, category, region). The ontology covers 2{,}112 ingredients across 49 classes (4 levels, 39 leaves) and 7 named relations (Table~\ref{tab:relations}). Dense retrieval uses \texttt{multilingual-e5-large} (1024-dim) with Pinecone as the vector store \cite{wang2024multilinguale5}. BM25 uses Okapi BM25 over dish name and ingredient text.

\subsection{Systems}

Task~1 compares four systems: \textbf{BM25} (keyword matching 
baseline), \textbf{BM25+Expansion} (BM25 with flat synonym expansion from the ingredient KB, no hierarchy), \textbf{Dense} (dense retrieval without ontology), and \textbf{Dense+Ontology} 
(dense retrieval with ontology query expansion, constraint filtering, and hierarchy-aware similarity). 

Task~2 compares three systems (BM25, Dense, Dense+Ontology) on the same diverse candidate set, and additionally reports a component ablation study to isolate the contribution of each ontology signal.


This comparison follows a standard sparse--dense setup: BM25 provides the lexical baseline, dense retrieval the semantic baseline, and Dense+Ontology tests whether structured knowledge 
adds gains beyond both \cite{robertson2009probabilistic, 
karpukhin2020dense, thakur2021beir}.
% ──────────────────────────────────────────────────────────────
\subsection{Evaluation Protocol}
\label{ssec:eval_protocol}

\textbf{Dense retrieval index.}
Each dish is indexed as a single document: \texttt{dish\_name} (repeated 3$\times$ for name-match boosting) + \texttt{category} + all \texttt{ingredient\_names} (Vietnamese). The embedding model (\texttt{multilingual-e5-large}, 1024-dim) encodes documents with the \texttt{passage:} prefix and queries with the \texttt{query:} prefix, following the E5 protocol \cite{wang2024multilinguale5}.

\textbf{Evaluation metrics.}
Both tasks are evaluated using ranking metrics including Precision@k and NDCG, where NDCG follows the cumulative-gain formulation of J\"arvelin and Kek\"al\"ainen \cite{jarvelin2002cg}.

% ──────────────────────────────────────────────────────────────
\subsubsection{Task~1: Class-Based Retrieval Protocol}
\label{sssec:eval_task1}
\hfill\break\indent\textbf{Query construction.}
Queries are generated programmatically from 10 Vietnamese-language templates (e.g., \textit{``các món \{ing\}''}, \textit{``\{a\} nấu với \{b\}''}) with class and method slots filled by random sampling over all 24 leaf ingredient classes and 10 cooking methods (seed~42). This yields 200 queries (50 per type) with deterministic ground-truth labels computed by the \texttt{FoodOntology} API: a dish is positive iff it contains $\geq$1 ingredient from every positive class, zero ingredients from any negative class, and matches the required cooking method. Labels require no human annotation and are deterministic given the ontology.

\textbf{Human validation of rule labels.}
To validate label quality, two independent annotators judged a random sample of 500 (query, dish) pairs. Inter-annotator agreement is substantial (Cohen's $\kappa = 0.62$, exact agreement 83.4\%). Rule labels agree with human majority consensus at $\kappa = 0.50$ (F1\,=\,0.72, recall\,=\,0.81). Most disagreements occur in multi-class queries ($\kappa = 0.31$) where annotators are uncertain about ingredient class membership (e.g., whether ``bột năng'' belongs to Starch); cooking-method ($\kappa = 0.63$) and negation ($\kappa = 0.57$) queries show strong agreement.

Task~1 involves no tunable hyperparameters.

% ──────────────────────────────────────────────────────────────
\subsubsection{Task~2: Related-Dish Recommendation Protocol}
\label{sssec:eval_task2}
\hfill\break\indent\textbf{Candidate construction.}
Anchor dishes (200) are selected via stratified sampling across 25 dish categories. For each anchor, 20 candidates are drawn from four sources to ensure diversity: (1)~top-5 by IDF-weighted Jaccard (easy for Jaccard-based systems), (2)~5 from mid-range Jaccard (rank 10--20), (3)~5 same-category dishes with Jaccard $<$ 0.2 (hard for Jaccard, tests ontology), and (4)~5 random negatives. This yields ${\sim}$4{,}000 pairs. Metrics (P@5, NDCG@5, MRR@5) use a positive threshold of mean judge score $\geq$\,1.0.

\textbf{LLM judge protocol.}
All pairs are scored by three LLM judges (\textbf{Llama-3.1 8B}, \textbf{Gemma-2 9B}, \textbf{Mistral 7B}) on a 0/1/2 relatedness scale with the following prompt:

\begin{quote}
\small
\textit{``Rate how related these two Vietnamese dishes are for a `similar dishes' recommendation. Score: 2 = very related (same type, similar ingredients), 1 = somewhat related (some overlap), 0 = unrelated. Reply with ONLY one number. Dish 1: \{dish\_a\}. Dish 2: \{dish\_b\}. Score:''}
\end{quote}

The three judges achieve Fleiss' $\kappa = 0.336$ (fair agreement) and pairwise exact agreement of 70--76\% (Llama--Gemma: 75.6\%, Llama--Mistral: 69.9\%, Gemma--Mistral: 74.1\%). Mean scores are consistent across judges (Llama: 0.79, Gemma: 0.79, Mistral: 0.97). The ground-truth label for each pair is the mean score across the three judges \cite{zheng2023judge, chen2024judgebias, huang2025llmjudge}.

\textbf{Human validation of LLM judges.}
Two independent annotators scored a random subset of 300 pairs on the same 0/1/2 scale. Inter-annotator agreement is moderate (Cohen's $\kappa_{\text{linear}} = 0.44$, exact agreement 60\%, adjacent agreement 96\%). Spearman correlation between the human consensus (mean of two annotators) and the LLM-judge mean is $\rho = 0.50$ ($p < 10^{-19}$), with Kendall $\tau = 0.44$. LLM judges exhibit a positive bias of $+0.28$ on the 0--2 scale, tending to overestimate relatedness. However, binarizing at the positive threshold ($\geq 1$), LLM judges achieve 96.1\% recall of human-relevant pairs, indicating that the automatic labels rarely miss truly related dishes. This moderate correlation is consistent with prior findings on LLM-as-judge reliability for subjective similarity tasks \cite{zheng2023judge}.

\textbf{Weight optimization.}
Similarity weights ($\alpha, \beta, \gamma, \delta$ in Eq.~\ref{eq:sim}) are tuned via 5-fold cross-validation over the 200 anchors: for each fold, weights are optimized on the training anchors (160) via Nelder-Mead to maximize Spearman correlation, and evaluated on the held-out anchors (40). Final weights are the mean across folds.

\subsection{Results}

\textbf{Task~1: Class-Based Retrieval} (Table~\ref{tab:task1}).
Dense+Ontology achieves P@20\,=\,0.446 and NDCG@20\,=\,0.472, outperforming Dense by +32\% and +37\% respectively, and BM25 by +94\% and +103\%. The gain from ontology hierarchy (+32\% over Dense) exceeds the gain from flat synonym expansion (+28\% BM25+Expansion over BM25), confirming that structured class expansion is more effective than flat lookup. Paired Wilcoxon signed-rank tests on per-query P@20 \cite{smucker2007sigtest} confirm all pairwise improvements are statistically significant ($p < 0.001$ for all pairs).

\begin{table}[t]
\centering
\caption{Task~1: Class-based retrieval (200 queries, top-20).}
\label{tab:task1}
\small
\begin{tabular}{lccc}
\toprule
\textbf{System} & \textbf{P@20} & \textbf{NDCG@20} & \textbf{MRR@20} \\
\midrule
BM25              & 0.230 & 0.232 & 0.397 \\
BM25+Expansion    & 0.295 & 0.298 & 0.428 \\
Dense             & 0.339 & 0.344 & 0.511 \\
\textbf{Dense+Ontology} & \textbf{0.446} & \textbf{0.472} & \textbf{0.711} \\
\bottomrule
\end{tabular}
\end{table}

\textbf{Task~2: Related-Dish Recommendation} (Tables~\ref{tab:task2_ablation} and~\ref{tab:task2}).

To avoid circularity from Jaccard-based candidate selection, we construct a diverse candidate set for each of 200 anchor dishes: 5 candidates from top Jaccard overlap, 5 from mid-range similarity, 5 from same-category but low Jaccard (testing ontology contribution), and 5 random negatives, yielding ${\sim}$4{,}000 pairs scored by three LLM judges.

We first determine the similarity formula via component ablation (Table~\ref{tab:task2_ablation}). Eight configurations are compared using 5-fold cross-validation, with weights optimized on each training fold via Nelder-Mead to maximize Spearman correlation. Each ontology component contributes measurably: adding ClassOverlap improves P@5 by +7.4\% over Jaccard alone, MethodMatch adds +2.9\%, and SemanticSim adds +0.5\% MRR@5. The best configuration (Full, all four components) yields optimized weights $\alpha = 0.40$, $\beta = 0.18$, $\gamma = 0.11$, $\delta = 0.31$, showing that ontology-derived components ($\beta + \gamma + \delta = 0.60$) account for 60\% of the similarity signal.


Using these optimized weights, Table~\ref{tab:task2} compares Dense+Ontology against BM25 and Dense baselines on the same diverse candidate set. Dense+Ontology outperforms both baselines across all metrics.

\begin{table}[t]
\centering
\caption{Task~2: Ablation study (5-fold CV, 200 anchors, diverse candidates). Positive threshold: mean judge score $\geq$ 1.0.}
\label{tab:task2_ablation}
\small
\begin{tabular}{lccc}
\toprule
\textbf{Configuration} & \textbf{P@5} & \textbf{NDCG@5} & \textbf{MRR@5} \\
\midrule
A: Jaccard only              & 0.741{\scriptsize±.061} & 0.755{\scriptsize±.053} & 0.855{\scriptsize±.042} \\
B: +ClassOverlap             & 0.796{\scriptsize±.051} & 0.816{\scriptsize±.038} & 0.905{\scriptsize±.028} \\
C: +MethodMatch              & 0.819{\scriptsize±.054} & 0.844{\scriptsize±.044} & 0.944{\scriptsize±.032} \\
\textbf{D: Full (all 4)}     & \textbf{0.819}{\scriptsize±.054} & \textbf{0.845}{\scriptsize±.043} & \textbf{0.949}{\scriptsize±.029} \\
\midrule
E: No Jaccard                & 0.815{\scriptsize±.051} & 0.839{\scriptsize±.047} & 0.939{\scriptsize±.045} \\
F: No ClassOverlap           & 0.811{\scriptsize±.055} & 0.831{\scriptsize±.047} & 0.923{\scriptsize±.037} \\
G: No MethodMatch            & 0.796{\scriptsize±.051} & 0.816{\scriptsize±.038} & 0.905{\scriptsize±.028} \\
H: No SemanticSim            & 0.819{\scriptsize±.054} & 0.844{\scriptsize±.044} & 0.944{\scriptsize±.032} \\
\bottomrule
\end{tabular}
\end{table}

\begin{table}[t]
\centering
\caption{Task~2: System comparison using optimized weights (200 anchors, diverse candidates).}
\label{tab:task2}
\begin{tabular}{lccc}
\toprule
\textbf{System} & \textbf{P@5} & \textbf{NDCG@5} & \textbf{MRR@5} \\
\midrule
BM25            & 0.784 & 0.812 & 0.913 \\
Dense           & 0.814 & 0.834 & 0.930 \\
\textbf{Dense+Ontology}  & \textbf{0.827} & \textbf{0.852} & \textbf{0.956} \\
\bottomrule
\end{tabular}
\end{table}

Across both tasks, ontology-enhanced retrieval outperforms baselines. For Task~1, Dense+Ontology improves NDCG@20 by +37\% and MRR@20 by +39\% over Dense, with hierarchy expansion contributing more than flat synonym expansion (+28\%).

For Task~2, Dense+Ontology outperforms both BM25 and Dense on the diverse candidate set (Table~\ref{tab:task2}). The ablation study (Table~\ref{tab:task2_ablation}) demonstrates clear incremental gains from each ontology component. Starting from Jaccard-only (P@5\,=\,0.741), adding ClassOverlap yields +7.4\%, MethodMatch adds +2.9\%, confirming that structured class knowledge and cooking-method matching complement ingredient overlap. The ``No Jaccard'' configuration (E) still achieves P@5\,=\,0.815, showing that ontology signals alone provide strong dish similarity even without direct ingredient matching.
